\subsection{严格态射的刻画}

命题 1. --- 设 E 与 F 是两个局部凸空间，u 是从 E 到 F 中的一个连续线性映射。为使 u 是严格态射，必要且充分的是满足下列条件：

(MS) 对每个半范数 \(p \in \mathrm{S}(\mathrm{E})\)（在 \(u\) 的核上为零的），存在 \(q\)（\(\mathrm{S}(\mathrm{F})\) 中的）使得 \(p \leqslant q \circ u\)。

设 N 是 u 的核，M 是 u 的像；我们引入 u 的典范分解，即

\[
\mathrm{E} \xrightarrow {\pi} \mathrm{E} / \mathrm{N} \xrightarrow {\tilde {u}} \mathrm{M} \xrightarrow {\iota} \mathrm{F}.
\]

E 上在 N 上为零的连续半范数是半范数 \(p_{1} \circ \pi\)（其中 \(p_{1}\) 取遍 S(E/N)）；类似地，S(M) 由这样的半范数 \(q_{1}\) 组成：存在 \(q \in \mathrm{S}(F)\) 使得 \(q_{1} \leqslant q/F\)。最后，u 是严格态射当且仅当双射连续线性映射 \(\tilde{u}\) 有连续的逆；这也就是说 S(E/N) 中的每个半范数都具有形式 \(q_{1} \circ \tilde{u}\)（其中 \(q_{1}\) 属于 S(M)）。命题 1 立即由这些注记得出。

命题 2. --- 设 E 与 F 是两个 Hausdorff 局部凸空间，u 是从 E 到 F 中的一个连续线性映射。为使 u 是严格态射，必要且充分的是它的转置 \({}^t u\): F' → E' 满足下列条件：

\begin{enumerate}
\def\labelenumi{\alph{enumi})}
\item
  \(^{t}u\) 的像在 \(E'\) 中对于 \(\sigma(E', E)\) 闭。
\item
  \(\mathbf{E}'\) 的每个等度连续子集，若包含在 \({}^t u\) 的像中，就是 \({}^t u\) 对 \(\mathbf{F}'\) 的某个等度连续子集所取的像。
\end{enumerate}

若如此，我们有 \(\operatorname{Ker}^t u = (\operatorname{Im} u)^\circ\) 与 \(\operatorname{Im}^t u = (\operatorname{Ker} u)^\circ\)，且存在从 \(\operatorname{Coker}^t u\) 到 \(\operatorname{Ker} u\) 的对偶上的以及从 \(\operatorname{Ker}^t u\) 到 \(\operatorname{Coker} u\) 的对偶上的典范同构。

设 N 是 u 的核，I 是 u 的像。由 II 第 47 页的推论 2，\(^{t}u\) 的核是 I 的正交，而 \(^{t}u\) 的像对于 \(\sigma(E', E)\) 的闭包是 N 的正交 \(N^{\circ}\)。于是 a) 与 b) 的合取等价于下列条件：

b') \(E'\) 的每个包含在 \(N^{\circ}\) 中的等度连续子集都是 \(^{t}u\) 对 \(F'\) 的某个等度连续子集所取的像。

由于 \(N^{\circ}\) 可以等同于 E/N 的对偶，IV 第 8 页的命题 9 (i) 表明：\(E'\) 的包含在 \(N^{\circ}\) 中的等度连续子集就是包含在形如 \(H_{p}\) 的集中的那些集，其中 p 是 E 上的连续半范数且在 N 上为零。于是条件 \(b')\) 说的是：对每个在 N 上为零的半范数 \(p \in \mathrm{S}(\mathrm{E})\)，存在 \(q \in \mathrm{S}(\mathrm{F})\) 使得 \(\mathrm{H}_{p} \subset {}^{t}u(\mathrm{H}_{q})\)。由 Hahn-Banach 定理（II，第 23 页，推论 1 与第 63 页，定理 1 与推论 1），我们有 \(u(\mathrm{H}_{q}) = \mathrm{H}_{q \circ u}\)，且关系 \(H_{p} \subset H_{p'}\) 与 \(p \leqslant p'\) 对所有半范数 p 与 \(p'\)（S(E) 中的）等价。因此关系 \(H_{p} \subset {}^{t}u(\mathrm{H}_{q})\) 等价于关系 \(p \leqslant q \circ u\)。由命题 1，性质 \(b')\) 蕴含 u 是严格态射。

假定 \(u\) 是严格态射。我们已看到，\(^t u\) 的核是 I 的正交，而 \(^t u\) 的像是 N 的正交。u 的余核是空间 F/I，它的对偶可以等同于 \(I^{\circ} = Ker^{t}u\)。类似地，N = Ker u 的对偶可以等同于 \(E^{\prime}/N^{\circ}\)（IV，第 8 页），即等同于 \(^{t}u\) 的余核，因为 \(N^{\circ}\) 是 \(^{t}u\) 的像。

注. --- 采用命题 2 的记号，性质 \(b'\) ) 还蕴含 \(u\) 对于弱化拓扑是严格态射（II，第 49 页，推论 3）。

命题 3. --- 设 E 与 F 是两个局部凸空间，u 是从 E 到 F 中的一个连续线性映射。我们假定 E 是 Hausdorff 的且 F 是可度量化的。为使 u 是严格态射，必要且充分的是 \({}^t u\) 的像在 E' 中对于弱拓扑 \(\sigma(E', E)\) 闭。

必要性由命题 2 得出。

反之，假定 \(^t u\) 的像对于 \(\sigma(E', E)\) 闭，并如命题 1 的证明中那样引入 \(u\) 的典范分解。由上述注记，逆映射 \(\tilde{u}^{-1}\)（\(\tilde{u}\) 的）对于弱化拓扑连续。但 F 的子空间 \(M = u(E)\) 是可度量化的，从而是有界型的（III，第 12 页，命题 2）；于是（IV，第 7 页，命题 7 (ii)），\(\tilde{u}^{-1}\) 连续，故 \(u\) 是严格态射。
% semantic-fragments: legacy-input-seam
\relax{}
